\chapter{1977年云南卷}

\section{解答题}

\begin{problem}
\begin{enumerate}
\item 计算 \(\left(\frac{8x^{-3}}{\sqrt{y^3}z^{-6}}\right)^{-\frac{1}{3}}\)（\(=\)？）

\item 计算 \(\cos\frac{3\pi}{8}\cdot\cos\frac{\pi}{8}\) 之值.

\item 计算 \(\log_3\frac{1}{27}+\log_{\frac{1}{2}}8+\log_2 8+\log_4 64\)（\(=\)？）

\item 解方程 \(2\sqrt{x-3}+6=x\).
\end{enumerate}
\end{problem}

\begin{solution}
\begin{enumerate}
\item 原式在实数范围内有意义时，必须有 \(x\neq0\)、\(y>0\) 且 \(z\neq0\). 由于 \(y>0\)，有 \(\sqrt{y^3}=y^{\frac{3}{2}}\)，于是
\[
\frac{8x^{-3}}{\sqrt{y^3}z^{-6}}
=\frac{2^3x^{-3}}{y^{\frac{3}{2}}z^{-6}}
=2^3x^{-3}y^{-\frac{3}{2}}z^6
\]
因此
\[
\left(\frac{8x^{-3}}{\sqrt{y^3}z^{-6}}\right)^{-\frac{1}{3}}
=\left(2^3x^{-3}y^{-\frac{3}{2}}z^6\right)^{-\frac{1}{3}}
=2^{-1}xy^{\frac{1}{2}}z^{-2}
=\frac{x\sqrt{y}}{2z^2}
\]

\item 由余函数关系，
\[
\cos\frac{3\pi}{8}
=\sin\left(\frac{\pi}{2}-\frac{3\pi}{8}\right)
=\sin\frac{\pi}{8}
\]
所以
\[
\cos\frac{3\pi}{8}\cdot\cos\frac{\pi}{8}
=\sin\frac{\pi}{8}\cos\frac{\pi}{8}
=\frac{1}{2}\sin\frac{\pi}{4}
=\frac{\sqrt{2}}{4}
\]

\item 将各真数写成相应底数的幂，得
\(\log_3\frac{1}{27}=\log_3 3^{-3}=-3\)，\(\log_{\frac{1}{2}}8=\log_{\frac{1}{2}}\left(\frac{1}{2}\right)^{-3}=-3\)
\(\log_2 8=\log_2 2^3=3\)，\(\log_4 64=\log_4 4^3=3\).
因此原式的值为 \(-3-3+3+3=0\).

\item 原方程有意义的条件是 \(x-3\geqslant0\). 令
\[
y=\sqrt{x-3}
\]
则 \(y\geqslant0\)，且 \(x=y^2+3\). 代入原方程，得
\[
2y+6=y^2+3
\]
即
\(y^2-2y-3=0\)，\(\left(y-3\right)\left(y+1\right)=0\).
所以 \(y=3\) 或 \(y=-1\). 由 \(y\geqslant0\)，只能取 \(y=3\)，从而 \(x=3^2+3=12\). 代回原方程，\(2\sqrt{12-3}+6=2\times3+6=12\)，故原方程的解为 \(x=12\).
\end{enumerate}
\end{solution}

\begin{problem}
某拖拉机厂为了支援农业生产，在预定期限内计划生产一批新型拖拉机，若按原来每天生产 \(20\text{~台}\)，则差 \(100\text{~台}\) 不能按时完成任务；在揭批“四人帮”的伟大运动的推动下，每天生产了 \(25\text{~台}\)，结果到期比原计划超额 \(50\text{~台}\)，问原计划生产拖拉机多少台？预定期限是多少天？
\end{problem}

\begin{solution}
设预定期限为 \(x\text{~天}\)，原计划生产拖拉机 \(y\text{~台}\). 按原计划每天生产 \(20\text{~台}\)，到期还差 \(100\text{~台}\)，所以
\[
y=20x+100
\]
每天生产 \(25\text{~台}\) 时，到期比原计划多生产 \(50\text{~台}\)，所以
\[
25x=y+50
\]
即 \(y=25x-50\). 因此根据题意，得方程组
\[
\begin{cases}
y=20x+100,\\
y=25x-50.
\end{cases}
\]
比较两式，得 \(20x+100=25x-50\)，所以 \(5x=150\)，从而 \(x=30\). 代入 \(y=20x+100\)，得
\[
y=20\times30+100=700
\]
由于 \(x>0\)、\(y>0\)，所得结果符合题意. 答：原计划生产拖拉机 \(700\text{~台}\)，预定期限是 \(30\text{~天}\).
\end{solution}

\begin{problem}
如图，设两圆相切于 \(C\) 点，\(PQ\) 是两圆的外公切线，\(A\)，\(B\) 分别为 \(PQ\) 与两圆相切的切点，试证明 \(AC\perp BC\).

\bitmapfigure[max width=6.8cm]{img/1977/yunnan/q3_fig1.png}
\end{problem}

\begin{solution}
过两圆的切点 \(C\) 作两圆的内公切线 \(CD\)，交外公切线 \(PQ\) 于 \(D\). 由同一点向同一圆所作的两条切线长相等，得
\(DA=DC\)，\(DB=DC\).
因此，在等腰三角形 \(ADC\) 和 \(BDC\) 中，分别有
\(\angle DAC=\angle DCA\)，\(\angle DCB=\angle DBC\).
又因为 \(A\)，\(D\)，\(B\) 共线，射线 \(DA\) 与 \(DB\) 方向相反，所以
\[
\angle ADC+\angle CDB=180^\circ
\]
将三角形 \(ADC\) 和 \(BDC\) 的内角和相加，得到
\[
\angle DAC+\angle DCA+\angle DCB+\angle DBC=180^\circ
\]
而射线 \(CD\) 位于 \(\angle ACB\) 的内部，故
\[
\angle ACB=\angle DCA+\angle DCB=90^\circ
\]
所以 \(AC\perp BC\).
\end{solution}

\begin{problem}
求过 \(x^2+y^2=4\) 上一点 \(P\left(1,\sqrt{3}\right)\) 的切线 \(L\) 的方程.
\end{problem}

\begin{solution}
圆 \(x^2+y^2=4\) 的圆心为 \(O\left(0,0\right)\)，半径为 \(2\). 先验证
\[
1^2+\left(\sqrt{3}\right)^2=4
\]
所以 \(P\left(1,\sqrt{3}\right)\) 在圆上. 半径 \(OP\) 垂直于过 \(P\) 的切线，故切线的法向量可取 \(\left(1,\sqrt{3}\right)\). 由点法式，切线 \(L\) 满足
\[
1\left(x-1\right)+\sqrt{3}\left(y-\sqrt{3}\right)=0
\]
整理得
\[
x+\sqrt{3}y-4=0
\]
代入点 \(P\) 可得 \(1+3-4=0\)，故该直线确实经过 \(P\). 答：切线 \(L\) 的方程为 \(x+\sqrt{3}y-4=0\).
\end{solution}

\begin{problem}
一船以 \(24\text{~浬/小时}\) 的速度向正北航行，在 \(A_1\) 点望见一灯塔 \(S\)，在船的北偏东 \(30^\circ\)，\(15\text{~分钟}\) 后，在 \(A_2\) 点望见灯塔在船的北偏东 \(65^\circ\)（如图），求船在 \(A_2\) 点时，与灯塔的距离 \(A_2S\).（已知 \(\sin35^\circ=0.574\)）

\bitmapfigure[max width=6.0cm]{img/1977/yunnan/q5_fig1.png}
\end{problem}

\begin{solution}
船沿正北方向航行，因此 \(A_1A_2\) 沿正北方向. 在 \(A_1\) 点，射线 \(A_1S\) 为北偏东 \(30^\circ\)，所以
\[
\angle SA_1A_2=30^\circ
\]
在 \(A_2\) 点，射线 \(A_2A_1\) 指向正南，而射线 \(A_2S\) 为北偏东 \(65^\circ\)，所以
\[
\angle SA_2A_1=180^\circ-65^\circ=115^\circ
\]
从而
\[
\angle A_1SA_2=180^\circ-30^\circ-115^\circ=35^\circ
\]
船行驶了 \(15\) 分钟，即 \(\frac{15}{60}\) 小时，所以
\[
A_1A_2=24\times\frac{15}{60}=6\text{~浬}
\]
在三角形 \(A_1A_2S\) 中，由正弦定理，
\[
\frac{A_2S}{\sin\angle SA_1A_2}
=\frac{A_1A_2}{\sin\angle A_1SA_2}
\]
所以
\[
A_2S=\frac{6\sin30^\circ}{\sin35^\circ}
=\frac{6\times\frac{1}{2}}{0.574}
\approx5.23\text{~浬}
\]

答：船在 \(A_2\) 点时与灯塔的距离约等于 \(5.23\text{~浬}\).
\end{solution}

\begin{problem}
直角三角形中的一个锐角为 \(\beta\)，面积是12.

\begin{enumerate}
\item 求此三角形的外接圆的面积；

\item 当 \(\beta\) 取何值时，外接圆的面积最小.
\end{enumerate}

\bitmapfigure[max width=4.8cm]{img/1977/yunnan/q6_fig1.png}
\end{problem}

\begin{solution}
\begin{enumerate}
\item 设直角三角形为 \(ABC\)，其中 \(\angle C=90^\circ\)、\(\angle B=\beta\)，外接圆半径为 \(R\). 因为直角三角形的斜边是外接圆的直径，所以 \(AB=2R\). 又因为 \(0<\beta<90^\circ\)，在直角三角形中有
\(AC=2R\sin\beta\)，\(BC=2R\cos\beta\).
由三角形面积为 \(12\)，得
\[
\frac{1}{2}\cdot2R\sin\beta\cdot2R\cos\beta=12
\]
即
\(R^2\sin\left(2\beta\right)=12\)，\(R^2=\frac{12}{\sin\left(2\beta\right)}\).
所以外接圆的面积为
\[
\pi R^2=\frac{12\pi}{\sin\left(2\beta\right)}
\]
答：此三角形外接圆的面积为 \(\frac{12\pi}{\sin\left(2\beta\right)}\).

\item 因为 \(0<2\beta<180^\circ\)，所以 \(0<\sin\left(2\beta\right)\leqslant1\). 因而
\[
\frac{12\pi}{\sin\left(2\beta\right)}\geqslant12\pi
\]
当且仅当 \(\sin\left(2\beta\right)=1\) 时取等号，即 \(2\beta=90^\circ\)，所以 \(\beta=45^\circ\). 此时外接圆面积的最小值为 \(12\pi\). 答：当 \(\beta=45^\circ\) 时，外接圆的面积最小.

\end{enumerate}
\end{solution}

\begin{problem}
一个三角形三边的比为 \(4:5:6\)，面积为 \(27\sqrt{7}\,\mathrm{~cm}^2\)，求三边之长.
\end{problem}

\begin{solution}
设三边长分别为 \(a=4x\mathrm{~cm}\)、\(b=5x\mathrm{~cm}\)、\(c=6x\mathrm{~cm}\)，其中 \(x>0\). 半周长为
\[
s=\frac{a+b+c}{2}=\frac{4x+5x+6x}{2}\mathrm{~cm}=\frac{15x}{2}\mathrm{~cm}
\]
因此
\(s-a=\frac{7x}{2}\mathrm{~cm}\)，\(s-b=\frac{5x}{2}\mathrm{~cm}\)，\(s-c=\frac{3x}{2}\mathrm{~cm}\).
由海伦公式，三角形面积满足
\[
27\sqrt{7}\mathrm{~cm}^2
=\sqrt{s\left(s-a\right)\left(s-b\right)\left(s-c\right)}
=\frac{15\sqrt{7}}{4}x^2\mathrm{~cm}^2
\]
约去 \(\sqrt{7}\mathrm{~cm}^2\)，得 \(\frac{15}{4}x^2=27\)，所以
\[
x^2=\frac{36}{5}
\]
由于 \(x>0\)，故 \(x=\frac{6\sqrt{5}}{5}\). 于是
\(a=\frac{24\sqrt{5}}{5}\mathrm{~cm}\)，\(b=6\sqrt{5}\mathrm{~cm}\)，\(c=\frac{36\sqrt{5}}{5}\mathrm{~cm}\).
三边满足三角形不等式，故结果符合题意. 答：三边之长分别为 \(\frac{24\sqrt{5}}{5}\mathrm{~cm}\)、\(6\sqrt{5}\mathrm{~cm}\)、\(\frac{36\sqrt{5}}{5}\mathrm{~cm}\).
\end{solution}

\section{参考题}

\begin{problem}
某生产队要建造一个体积为 \(50\text{~立方米}\) 的有盖圆柱形氨水池（如图），问这个氨水池高和底半径取多大时用料最省（即表面积最小）？

\bitmapfigure[max width=4.6cm]{img/1977/yunnan/q8_fig1.png}
\end{problem}

\begin{solution}
设底面半径为 \(r\text{~m}\)，高为 \(h\text{~m}\)，则 \(r>0\)、\(h>0\). 由圆柱体积为 \(50\text{~m}^3\)，得
\(\pi r^2h=50\)，\(h=\frac{50}{\pi r^2}\).
氨水池有盖，所以表面积为两个底面面积与侧面积之和，即
\[
S=2\pi r^2+2\pi rh
=2\pi r^2+\frac{100}{r}
\]

方法一：将 \(S\) 看作 \(r\) 的函数，则
\(S'(r)=4\pi r-\frac{100}{r^2}\)，\(S''(r)=4\pi+\frac{200}{r^3}>0\).
所以 \(S\) 在 \(r>0\) 上为严格凸函数. 令 \(S'(r)=0\)，得
\[
4\pi r-\frac{100}{r^2}=0
\Longleftrightarrow
4\pi r^3=100
\Longleftrightarrow
r=\sqrt[3]{\frac{25}{\pi}}
\]
这是正数范围内的唯一驻点，且由 \(S''(r)>0\) 可知此时表面积取得最小值. 代入体积条件，
\[
h=\frac{50}{\pi r^2}=2\sqrt[3]{\frac{25}{\pi}}=2r
\]

方法二：由上式，
\[
S=2\pi r^2+\frac{50}{r}+\frac{50}{r}
\geqslant3\sqrt[3]{\left(2\pi r^2\right)\cdot\frac{50}{r}\cdot\frac{50}{r}}
=3\sqrt[3]{5000\pi}
\]
当且仅当
\[
2\pi r^2=\frac{50}{r}=\frac{50}{r}
\]
时取等号，解得 \(r=\sqrt[3]{\frac{25}{\pi}}\). 此时 \(h=\frac{50}{\pi r^2}=2r\). 因而两种方法均得到：当底面半径为 \(\sqrt[3]{\frac{25}{\pi}}\text{~m}\)、高为 \(2\sqrt[3]{\frac{25}{\pi}}\text{~m}\)，即高等于底面直径时，用料最省. 答：底面半径取 \(\sqrt[3]{\frac{25}{\pi}}\text{~m}\)，高取 \(2\sqrt[3]{\frac{25}{\pi}}\text{~m}\).
\end{solution}

\begin{problem}
求 \(\int_0^2\frac{x}{\sqrt{1+x^2}}\,\mathrm{d}x\).
\end{problem}

\begin{solution}
令 \(u=1+x^2\)，则 \(\mathrm{d}u=2x\,\mathrm{d}x\). 当 \(x=0\) 时 \(u=1\)，当 \(x=2\) 时 \(u=5\). 因此
\[
\int_0^2\frac{x}{\sqrt{1+x^2}}\,\mathrm{d}x
=\frac{1}{2}\int_1^5u^{-\frac{1}{2}}\,\mathrm{d}u
=\left.\sqrt{u}\right|_1^5
=\sqrt{5}-1
\]
\end{solution}
